% !TeX root = ../tesis.tex

\chapter{Antecedentes y Marco Teórico}
\label{ch:marco-teorico}

Al inicio del siglo XX, la capital mexicana contaba con escasos tranvías
eléctricos y unos pocos transportes de ruta, pero a medida que avanzó el siglo,
la población también creció, lo que obligó al gobierno federal a invertir en los
transportes colectivos. Este avance se vio frenado con el estallido de la
Revolución Mexicana en 1910, movimiento que no cesó hasta mediados de 1930. Para
esos años ya se habían implementado nuevas rutas de transporte colectivo, como
la creación del naciente tranvía y sus rutas que partían del centro de la
Ciudad. Otro avance notorio fue la promulgación de nuevos transportes de ruta a
localidades alejadas, siendo estos mayormente camiones. Sin embargo, el interés
por la expansión de las áreas urbanas no se dio hasta la época llamada ``El
Milagro Mexicano'', la cual abarca desde los inicios de 1940 hasta mediados de
la década de 1970, según \citet{molina2018}.

Es a partir del Cardenismo que la reconstrucción de la economía, la
infraestructura y el tejido político se consumó, haciendo necesaria una
inversión significativa sobre las ciudades y su comunicación. En su informe
entregado ante las cámaras federales de 1935, el Ejecutivo reportó una
recuperación notoria en las arcas públicas, destinando buena parte del cálculo
presupuestal a las necesidades sociales y de comunicación:

\begin{quote}
``Es verdaderamente satisfactorio para el Encargado del Poder Ejecutivo de la
Nación dirigirse a sus conciudadanos con motivo del primero de año, enviándoles
un mensaje cordial y lleno de optimismo, pues las condiciones del país nos
permiten afirmar que algunos de los anhelos trascendentales del pueblo podrán
realizarse totalmente en el presente ciclo, e iniciarse la ejecución de otros
varios que complementen la base de un período nuevo de positivo progreso.''
\citep[pp.~2--4]{biblioteca1935}
\end{quote}

Con este contexto, el urbanismo experimentó un acelerado crecimiento, el mismo
que se empalma con las necesidades de un país que transitaba de los caudillos y
lo rural hacia lo institucional y urbano.

% -----------------------------------------------------------------------
\section{La Hegemonía del Automóvil y la adopción de la automovilidad}
\label{sec:hegemonia-automovil}
% -----------------------------------------------------------------------

El problema de movilidad que actualmente afecta a la \cdmx{} y su periferia no
es un fenómeno coyuntural, sino el resultado de una adopción histórica y
cultural de la automovilidad como paradigma urbano. Este concepto, definido por
\citet{franco2022} en su tesis sobre el periodo 1903--1933, establece que el
vehículo privado dejó de ser un simple medio de transporte para convertirse en
el eje estructurador de la política pública y el imaginario social.

Esta adopción se consolidó a mediados del siglo XX con la priorización política
de la infraestructura vial sobre otros modos de transporte. El crecimiento
vehicular en el Distrito Federal fue exponencial, pasando de 150,584 unidades en
1955 a 209,904 en 1958 \citep{perlo2023}, una demanda que el Estado decidió
gestionar no mediante la diversificación modal, sino a través de la expansión de
la red de carreteras urbanas, sellando así la dependencia estructural del
automóvil que caracteriza la problemática actual.

